% ===================================================================
%  ТАБЛИЦА № 1 — Школа. — Лицей. — Класс.
%
%  Ceci n'est PAS une transcription : c'est une traduction, faite en
%  2026, en russe d'aujourd'hui. D'ou trois differences avec texto/io
%  et texto/fr, et trois seulement :
%
%   * pas de \nl ni de \cc — la traduction ne suit aucune ligne
%     imprimee, et les lignes de ce fichier ne sont que du confort ;
%   * pas de \begin{VUpage} — elle n'a ni feuillet ni folio, et la
%     page de lecture ne lui en affiche pas ;
%   * le gras se pose sur le TERME seul, ponctuation dehors.
%
%  Tout le reste est identique, et doit l'etre : les cles « %%K » sont
%  celles de texto/io, macro pour macro pour l'apparat, et les renvois
%  \textsuperscript{(n)} visent les MEMES objets de la planche — ce
%  sont eux qui ouvrent les gros plans.
%
%  L'ido est le texte de base ; le francais sert de controle, et
%  l'anglais et l'espagnol de calibrage : les six traductions doivent
%  s'accorder entre elles autant qu'avec la source.
%
%  UNE LANGUE A CAS. Le renvoi suit son substantif, et le substantif
%  se decline : le gras porte donc la forme FLECHIE, celle que la
%  phrase demande — « на \VUgras{скамьях} », non le nominatif. C'est
%  le nom de l'objet tel qu'il se lit, et objekti.py le releve tel
%  quel ; le gros plan s'en trouve intitule comme la phrase le dit.
%
%  LE NOM DE L'EDITEUR SE PRONONCE « del-MASS ». Delmas est un nom du
%  Midi -- « del mas », celui du mas -- et son s final se fait
%  entendre, comme dans Chaban-Delmas ; l'imprimerie est d'ailleurs a
%  Bordeaux, rue Saint-Christoly. La transcription le disait « Delma »
%  et perdait la syllabe. Le russe, qui decline les noms d'homme
%  termines par une consonne, demande alors le genitif : « таблицам
%  Дельмаса ».
% ===================================================================

%%K t01-apar-0 apar
\VUsaut{2.0mm}
\VUcentre{12.6pt}{120}{ПОЯСНИТЕЛЬНАЯ КНИЖКА}
\VUsaut{3.2mm}
\VUcentre{8.4pt}{160}{К}
\VUsaut{2.6mm}
\VUtitre{15.6pt}{0}{вспомогательным таблицам Дельмаса}
\VUsaut{3.4mm}
\VUornamento{9pt}
\VUsaut{5.0mm}
\VUcentre{13.2pt}{90}{ПЕРВАЯ СЕРИЯ}
\VUsaut{3.0mm}
\VUfilet{16mm}
\VUsaut{5.6mm}

%%K t01-apar-1 apar
\VUtitre{15.0pt}{60}{ТАБЛИЦА № 1}
\VUsaut{1.4mm}
\VUtitre{11.4pt}{0}{Школа. --- Лицей. --- Класс.}
\VUsaut{2.2mm}
\VUfilet{20mm}
\VUsaut{6.4mm}

%%K t01-tit-1 sub
\VUpk{10.2pt}{40}{Общее описание.}
\VUsaut{3.0mm}

%%K t01-01-1 p
1. --- На этой таблице изображён \VUgras{класс} в \VUgras{лицее}. В этом
классе восемь \VUgras{столов} \textsuperscript{(18)} и восемь \VUgras{скамей}
\textsuperscript{(32)}. На каждой скамье сидят три ученика. \VUgras{Учитель}
\textsuperscript{(7)} сидит на \VUgras{стуле} \textsuperscript{(8)} за своей
\VUgras{кафедрой} \textsuperscript{(6)}.

%%K t01-02-1 p
2. --- Слева от кафедры стоит застеклённый \VUgras{шкаф} \textsuperscript{(15)}.
Справа --- \VUgras{классная доска} \textsuperscript{(1)} с \VUgras{тряпкой}
\textsuperscript{(2)} и \VUgras{мелом} \textsuperscript{(3)}.

%%K t01-02-2 p
Над кафедрой висят \VUgras{настенные таблицы} \textsuperscript{(9, 11, 12)} и
\VUgras{карта} \textsuperscript{(10)}.

%%K t01-03-1 p
3. --- Не все ученики на своих \VUgras{местах} \textsuperscript{(17)}. Иван
\textsuperscript{(4)} стоит у доски; он держит тряпку и стирает
\VUgras{геометрические фигуры}.

%%K t01-03-2 p
Пётр стоит у кафедры; он собирает \VUgras{письменные работы}
\textsuperscript{(5)} и несёт их учителю.

%%K t01-04-1 p
4. --- \VUgras{Этот} учитель --- преподаватель \VUgras{новых языков}. Он
учит учеников говорить, читать и писать на иностранных языках
\VUgras{(немецком, английском, испанском, итальянском, русском} или
\VUgras{французском)}.

%%K t01-05-1 p
5. --- Класс очень большой и очень светлый. В глубине широкое
\VUgras{окно} с двумя \VUgras{фрамугами} \textsuperscript{(78)} и \VUgras{застеклённая
дверь}, чтобы его проветривать и освещать.

%%K t01-05-2 p
В этом классе не холодно. Посередине, чтобы его обогревать, стоит очень
хорошая \VUgras{печь} \textsuperscript{(46)} со своей \VUgras{трубой}
\textsuperscript{(45)}.

%%K t01-05-3 p
К перегородке прибиты \VUgras{вешалки} \textsuperscript{(58)}; на них висят
шляпы и пальто учеников.

%%K t01-tit-2 sub
\VUsaut{5.2mm}
\VUpk{10.2pt}{40}{Двор.}
\VUsaut{3.6mm}

%%K t01-06-1 p
6. --- \VUgras{Часы} \textsuperscript{(63)} показывают пять минут десятого.

%%K t01-06-2 p
\VUgras{Перемена} \textsuperscript{(76)} только что кончилась, и урок начинается
снова. Перемена проходит во \VUgras{дворе} \textsuperscript{(75)}. Мы видим
нескольких играющих учеников. \VUgras{Надзиратель} \textsuperscript{(74)} следит
за ними.

%%K t01-07-1 p
7. --- Во дворе мы видим и других людей: \VUgras{инспектора}
\textsuperscript{(73)}, \VUgras{директора} \textsuperscript{(71)} и \VUgras{завуча}
\textsuperscript{(72)}.

%%K t01-07-2 p
Инспектор проверяет школы, директор управляет заведением, завуч следит
за работой учеников.

%%K t01-07-3 p
Четвёртый --- \VUgras{сторож} \textsuperscript{(77)}. Он несёт под мышкой журнал,
чтобы отметить отсутствующих.

%%K t01-08-1 p
8. --- В глубине двора находятся \VUgras{гимнастические снаряды}
\textsuperscript{(70)} и мастерская \VUgras{ручного труда} \textsuperscript{(69)}.

%%K t01-08-2 p
Справа на таблице едва видны \VUgras{классы} \VUgras{музыки} и
\VUgras{рисования}.

%%K t01-noto-1 noto
\VUnotes{143.2mm}{%
(*) Для \textit{имён} тоже советуем передавать их в латинской форме.
Только так можно избежать бесчисленных разночтений. Да и как выбрать
между \textit{Giovanni, Jean, Johann, Jan, Hans, John, Иван} и прочими?
Не создавать же особый словарь имён? Возьмём латинский именительный
падеж и скажем: \VUgras{Ioannes, Ioanna; Iulius, Iulia; Iakobus;
Andreas; Lukas.} (Полная грамматика).%
}

%%K t01-tit-3 sub
\VUpk{10.2pt}{40}{Подробное описание.}
\VUsaut{2.4mm}

%%K t01-tit-4 sub
\VUpk{10.2pt}{40}{Хорошие ученики.}
\VUsaut{3.0mm}

%%K t01-09-1 p
9. --- Ученики на первой скамье справа от кафедры --- \VUgras{Генрих}
(*), \VUgras{Стефан} и \VUgras{Карл}.

%%K t01-09-2 p
Генрих очень внимателен и прилежен; он слушает учителя. На столе у него
\VUgras{тетрадь} \textsuperscript{(28)}, \VUgras{книга} \textsuperscript{(29)},
\VUgras{словарь} \textsuperscript{(30)} и \VUgras{пенал} \textsuperscript{(31)}. В его книге
много \VUgras{страниц}; в каждой главе несколько \VUgras{абзацев}, а в
каждой \VUgras{строке} много \VUgras{слов}.

%%K t01-09-4 p
Его сосед Стефан \textsuperscript{(24)} усерден. Он записывает в тетрадь урок,
заданный на следующий раз. У него прекрасный \VUgras{почерк}. Книга его
закрыта. Виден только \VUgras{переплёт} \textsuperscript{(27)}. Рядом с ним лежит
его \VUgras{циркуль} \textsuperscript{(26)}.

%%K t01-09-5 p
Карл \textsuperscript{(23)} держит книгу в руке; он открывает её, чтобы ещё раз
повторить урок. Как у всякого ученика, перед ним стоит
\VUgras{чернильница} \textsuperscript{(25)}. Он послушен.

%%K t01-10-1 p
10. --- За следующим столом сидят \VUgras{Людовик, Антон} и \VUgras{Павел}.
Людовик отвечает свой \VUgras{урок} \textsuperscript{(22)} и отвечает на
\VUgras{вопросы} учителя. Антон \textsuperscript{(21)} очень невнимателен; иногда
учитель даже наказывает его. Павел \textsuperscript{(20)} --- хороший товарищ,
приветливый и услужливый. Он очень внимателен.

%%K t01-tit-5 sub
\VUsaut{4.2mm}
\VUpk{10.2pt}{40}{Плохие ученики.}
\VUsaut{3.2mm}

%%K t01-11-1 p
11. --- Позади Генриха сидит \VUgras{Георгий} \textsuperscript{(38)}. Он не злой
мальчик, но \VUgras{болтлив}, невнимателен и непоседлив. Его
\VUgras{легкомыслие} и \VUgras{непослушание} вносят \VUgras{беспорядок} в
урок, и он часто заслуживает строгого \VUgras{замечания}. Его
\VUgras{черновик} \textsuperscript{(35)} грязен, он ставит на нём \VUgras{кляксы}
своим \VUgras{пером} \textsuperscript{(34)}, отчего постоянно пускает в ход
\VUgras{резинку} \textsuperscript{(36)}; его \VUgras{ранец} \textsuperscript{(33)} разорван,
до того он неаккуратен. Зачем он приносит свою \VUgras{ручку}
\textsuperscript{(37)} в класс новых языков?

%%K t01-11-2 p
Его сосед Эдмунд \textsuperscript{(39)} его не любит. Правда, он немного ленив,
но тих и спокоен. Он не болтает; только вместо того, чтобы слушать,
предпочитает читать \VUgras{рассказы} в своей \VUgras{хрестоматии}.

%%K t01-11-3 p
Третий ученик на этой скамье --- \VUgras{Людвиг} \textsuperscript{(40)}. Он
старателен. Он пишет на белом \VUgras{листе бумаги}. Перед собой он
положил \VUgras{промокательную бумагу} \textsuperscript{(41)}.

%%K t01-12-1 p
12. --- Ученики на второй скамье, слева от учителя, совсем малы. Первый,
маленький \VUgras{Эмиль}, ещё не умеет хорошо писать; он приносит свою
\VUgras{грифельную доску} \textsuperscript{(42)}, свой \VUgras{грифель} и свою
\VUgras{губку} \textsuperscript{(43)}.

%%K t01-12-2 p
Рядом с ним \VUgras{Август} и \VUgras{Бернард} \textsuperscript{(44)}. Последний
хорошо работает, но \VUgras{вид} у него очень неопрятный.

%%K t01-13-1 p
13. --- Позади них сидят трое \VUgras{лентяев}. \VUgras{Альберт} не учит
уроков. \VUgras{Яков} \textsuperscript{(47)} спит вместо того, чтобы слушать. Его
\VUgras{лень} навлекает на него \VUgras{выговоры} и даже \VUgras{наказания}.
Наконец, \VUgras{Христиан} \textsuperscript{(48)} слишком легкомыслен. Он плохо
\VUgras{ведёт себя} и не старается исправиться.

%%K t01-14-1 p
14. --- Их соседи, напротив, ведут себя хорошо. \VUgras{Эрнест}
\textsuperscript{(51)} любит порядок и точность; он всегда пользуется своей
\VUgras{линейкой} \textsuperscript{(50)} и своим \VUgras{карандашом} \textsuperscript{(49)}.
Он \VUgras{приходящий ученик}.

%%K t01-14-2 p
\VUgras{Франциск} \textsuperscript{(52)} --- \VUgras{пансионер}; он носит форму, ест
и спит в заведении. Он хороший \VUgras{товарищ}.

%%K t01-14-3 p
Мориц чинит свой \VUgras{карандаш} \VUgras{перочинным ножом} \textsuperscript{(56)};
он очень бережлив.

%%K t01-14-4 p
Он приносит все свои книги, свою \VUgras{грамматику} \textsuperscript{(55)}, свою
\VUgras{записную книжку} \textsuperscript{(54)} и даже \VUgras{блокнот}
\textsuperscript{(53)}. Его \VUgras{портфель} \textsuperscript{(57)} лежит на полу позади
него.

%%K t01-14-5 p
Два стола в глубине класса заняты \VUgras{хорошими учениками}
\textsuperscript{(19)}.

%%K t01-tit-6 sub
\VUsaut{4.6mm}
\VUpk{10.2pt}{40}{Рисование и музыка.}
\VUsaut{3.4mm}

%%K t01-15-1 p
15. --- В \VUgras{классе рисования} на стене висят красивые \VUgras{образцы},
а с потолка свешивается \VUgras{лампа} \textsuperscript{(62)} с \VUgras{абажуром}.
\VUgras{Учитель} \textsuperscript{(60)} очень терпелив; он исправляет
\VUgras{рисунки} \textsuperscript{(61)} учеников и учит их рисовать \VUgras{бюст}
\textsuperscript{(59)} с натуры.

%%K t01-16-1 p
16. --- \VUgras{Класс музыки} кажется очень занятным. Там учатся читать
\VUgras{ноты}, петь \VUgras{ступени гаммы} \textsuperscript{(67)}, написанные на
\VUgras{нотном стане} (до, ре, ми, фа, соль, ля, си\,=\,C. D. E. F. G. A.
B.), знать \VUgras{ключи} и петь прелестные \VUgras{мелодии}. Ноты кладут
на \VUgras{пюпитр} \textsuperscript{(65)}. Один из учеников \VUgras{играет на
рояле} \textsuperscript{(64)}, а \VUgras{учитель музыки} \textsuperscript{(66)}
\VUgras{отбивает такт} \textsuperscript{(68)}.

%%K t01-17-1 p
17. --- Едва виден угол мастерской \VUgras{ручного труда} \textsuperscript{(69)},
где дети могут учиться строгать дерево и опиливать железо. Не в каждой
школе она есть.

%%K t01-tit-7 sub
\VUsaut{5.0mm}
\VUpk{10.2pt}{40}{Класс естественных наук.}
\VUsaut{3.6mm}

%%K t01-18-1 p
18. --- Учитель математики преподаёт ученикам \VUgras{геометрию,}
\VUgras{арифметику, счёт} и \VUgras{метрическую систему}. На таблице мы
видим главные геометрические фигуры: \VUgras{круг} \textsuperscript{(a)},
\VUgras{квадрат} \textsuperscript{(b)}, \VUgras{треугольник} \textsuperscript{(c)},
\VUgras{линию} \textsuperscript{(d)}.

%%K t01-18-2 p
Мы видим также \VUgras{действие}; это не \VUgras{сложение}, не
\VUgras{вычитание} и не \VUgras{умножение}: это \VUgras{деление}
\textsuperscript{(e)}.

%%K t01-19-1 p
19. --- На таблице метрической системы мы видим: \VUgras{метр}
\textsuperscript{(a)}, \VUgras{весы} \textsuperscript{(b)} и их \VUgras{гири},
\VUgras{литр} \textsuperscript{(e)}, \VUgras{декалитр} \textsuperscript{(f)},
\VUgras{децилитр} и \VUgras{кубический метр} \textsuperscript{(d)}.

%%K t01-19-2 p
Ученики изучают также \VUgras{естественные науки} \textsuperscript{(11)} ---
зоологию \textsuperscript{(a)}, ботанику \textsuperscript{(b)}, геологию
\textsuperscript{(c)} --- и \VUgras{историю} \textsuperscript{(12)}.

%%K t01-19-3 p
В шкафу стоят приборы для \VUgras{физики} \textsuperscript{(15)} и для
\VUgras{химии} \textsuperscript{(16)}.

%%K t01-19-4 p
На шкафу стоят: \VUgras{шар} \textsuperscript{(14)}, \VUgras{конус} и
\VUgras{глобус} \textsuperscript{(13)}.

%%K t01-19-5 p
Над таблицами висит \VUgras{карта}, изображающая пять частей света:
\VUgras{Европу} \textsuperscript{(g)}, \VUgras{Азию} \textsuperscript{(e)},
\VUgras{Африку} \textsuperscript{(f)}, \VUgras{Америку} \textsuperscript{(ab)} и
\VUgras{Океанию} \textsuperscript{(h)}, и два великих океана ---
\VUgras{Тихий} \textsuperscript{(c)} и \VUgras{Атлантический} \textsuperscript{(d)}.
\VUsaut{6.0mm}
\VUfilet{22mm}
